\section{Fundamentação teórica}

A segurança ofensiva de redes compreende um conjunto de técnicas e metodologias destinadas a identificar, explorar e relatar vulnerabilidades em sistemas e infraestruturas de rede, com o propósito proativo de fortalecer suas defesas. Diferente da segurança defensiva, que visa proteger sistemas contra ataques, a abordagem ofensiva simula as ações de adversários reais para avaliar a resiliência dos mecanismos de proteção. Como mostrado no trabalho de \citet{offsec-emulation}, a área da segurança vem ganhando atenção por se mostrar eficiente e ajudar no treinamento de equipes e refinamento nas técnicas de detecção de ataques.

Neste projeto, a segurança ofensiva é aplicada por meio de ferramentas que realizam varreduras, explorações e testes de estresse em protocolos de rede, tanto em ambientes cabeados quanto sem fio. Técnicas como flooding (inundação), deautenticação em redes 802.11, spoofing (falsificação) de endereços IP e MAC são fundamentais para o teste de penetração, explorando falhas em protocolos, configurações inadequadas ou limitações de implementação.

Para a implementação de uma ferramenta ofensiva eficaz, é essencial o domínio dos protocolos de rede que atuam em diversas camadas do modelo TCP/IP. O TCP (\textit{Transmission Control Protocol}), protocolo orientado a conexão da camada de transporte, é alvo de técnicas como varredura e flooding. O UDP (\textit{User Datagram Protocol}), não orientado a conexão, permite varredura de portas com payloads específicos, explorando sua natureza \textit{connectionless}. O ICMP (\textit{Internet Control Message Protocol}), utilizado para diagnóstico, serve de base para ataques de flooding e técnicas de reflexão. Para redes sem fio, o padrão IEEE 802.11 (Wi-Fi) é essencial, sendo explorado por meio da captura de beacons e ataques de desautenticação, que manipulam quadros de gerenciamento.

O ecossistema de segurança ofensiva conta com diversas ferramentas consolidadas, cada uma com seus pontos fortes e limitações. Ferramentas como Nmap, para varredura de rede, e Scapy, para manipulação de pacotes em Python, são amplamente utilizadas, mas podem apresentar limitações de desempenho ou modularidade. Suítes como Aircrack-ng focam em redes Wi-Fi, mas com integração limitada a ambientes cabeados. Observa-se uma lacuna para uma ferramenta unificada.

A linguagem de programação Rust surge como uma escolha estratégica para o desenvolvimento de ferramentas de segurança ofensiva devido a um conjunto singular de vantagens. Seu sistema de propriedade (\textit{ownership}) e empréstimo (\textit{borrowing}) garante segurança de memória em tempo de compilação, prevenindo vulnerabilidades comuns como \textit{buffer overflows} e \textit{use-after-free}, problemas frequentes em linguagens como C/C++. Rust gera código máquina otimizado com desempenho similar a essas linguagens de baixo nível, crucial para operações de alta velocidade como flooding e varredura massiva. A linguagem também oferece concorrência segura, prevenindo data races em tempo de compilação, o que permite a implementação robusta de técnicas \textit{multithreading}. A interoperabilidade eficiente com C via FFI (\textit{Foreign Function Interface}) permite acesso direto a sockets brutos e chamadas de sistema. A combinação de desempenho, segurança e um ecossistema de desenvolvimento produtivo faz de Rust uma base ideal para construir ferramentas ofensivas confiáveis e eficientes.